\section{Conclusion}
\label{sec:conclusion}

This dissertation examined unusually large trades and short-horizon probability adjustment on Polymarket during the 2026 FIFA World Cup. The canonical sample comprises 6,725,732 Polymarket Data API trade records across 360 markets. Dune market-level counts served only as an external completeness check and did not contribute observations to the regression data.

Large trades are classified by market-specific, past-only P99 thresholds under expanding and rolling seven-day histories, each requiring 1,000 prior trades. Prices are aligned without interpolation using a pre-interval baseline and post-horizon outcome. Fixed-effects regressions, lead--lag checks, alternative binary-response models, prediction, and anomaly screening address distinct questions.

The principal finding is a robust but economically small directional association. The large-minus-ordinary five-minute contrasts are 0.000082 under expanding history and 0.000081 under rolling history. Under the primary expanding-P99 signed-log specification, a one-standard-deviation increase in large-trade flow is associated with approximately 0.0224 percentage points of probability movement, compared with 0.0059 percentage points for ordinary flow. Estimates at minus fifteen, minus five, and minus one minute are approximately zero, while positive associations appear at plus one, plus five, and plus fifteen minutes. The plus-thirty-minute result is not robust across both definitions. The association remains positive under event-by-minute and event-by-five-minute fixed-effect specifications. These controls reduce confounding from information common to related contracts within the same match-time cell, but they cannot remove contract-specific responses to contemporaneous match news; the estimates remain associational.

Mechanism evidence is qualified. Linear probability models suggest a larger large-flow association with update incidence, whereas market-fixed-effects logit differences are insignificant. H3a is supported in the unconditional minute-5-to-minute-15 specification, whereas the conditional non-zero-initial-movement evidence is not statistically significant; the result is an average subsequent-continuation association rather than evidence of strict persistence. Null Brier results provide no evidence of superior information. Supervised models predict reversal, not manipulation, while unsupervised anomalies are review candidates rather than confirmed misconduct.

These results support a deliberately narrow conclusion: unusually large directional flow is followed by a small short-horizon movement in the same direction, and this temporal pattern is robust to two past-only threshold definitions and demanding event-time controls. They do not establish that the trade caused the movement, that the trader possessed private information, or that anomalous behaviour constitutes manipulation. Wallet-level activity, cumulative gross trading volume, and individual trade size remain distinct constructs and are interpreted separately.

The dissertation makes three connected contributions. First, it replaces a
future-informed full-window size classification with expanding and rolling
thresholds that can be reproduced using only information available before each
classification day. Second, it combines strict price alignment with placebo
leads, post-trade lags, and event-by-time comparisons, making the temporal basis
of the primary association explicit. Third, it separates explanatory
regression, supervised prediction, and unsupervised anomaly discovery. This
separation prevents predictive performance or unusualness from being presented
as proof of a behavioural mechanism or integrity violation.

Supplementary wallet and concentration features improve reversal prediction,
but fixed-effects tests show no general push-and-reverse pattern. The contextual
audit reduces 73 rows to 35 linked or standalone units and establishes no
manipulation; it prioritises investigation rather than labelling traders.

The principal limitations are observational identification, the absence of independently timestamped public-news covariates, imperfect correspondence between wallet addresses and economic owners, dependence on an operational P99 threshold, and external validity from a single high-attention tournament. Even event-by-time fixed effects cannot absorb information specific to one contract within a match-time cell. Likewise, a public address may combine several strategies or represent only part of one economic actor's activity. These measurement limits make causal or identity-based claims inappropriate.

Future work could add timestamped match news and external probabilities,
replicate the design across event types, preregister thresholds, and test wallet
patterns over longer panels. The versioned framework keeps claims proportionate
to the evidence while documenting assumptions for replication.
